\section{Reduced-order thermoelastic analysis of the strip}
\label{sec:reduced}

This section derives the closed-form relations used to size the deformation before any coupled
computation. All of them follow from the model of Section~\ref{sec:formulation} under the stated
assumptions; the numerical values at the baseline are collected in Section~\ref{sec:baseline}.

\subsection{Free thermal curvature and the sign chain}
\label{sec:freecurv}
Let $\eta$ be the coordinate across the strip measured from the mid-plane (positive towards the outer
face), and let the temperature be linear,
$\theta(\eta)=\bar\theta+\Gamma\eta$, with $\Gamma=\Delta T_t/d>0$ by condition B5. For a thin strip with
traction-free faces $\sigma_{yy}\simeq0$, and plane strain gives
$\sigma_{xx}=E'[\varepsilon_{xx}-(1+\nu)\alpha\theta]$ with $E'=E/(1-\nu^2)$. With
$\varepsilon_{xx}=\varepsilon_0+\kappa\eta$, a vanishing axial force resultant gives
$\varepsilon_0=(1+\nu)\alpha\bar\theta$, and a vanishing moment resultant gives the free thermal
eigen-curvature
\begin{equation}
\kth=(1+\nu)\,\alpha\,\frac{\Delta T_t}{d}=(1+\nu)\,\alpha\,\frac{\qpp}{k_s},
\qquad \Delta T_t=\frac{\qpp d}{k_s}, \label{eq:kappa}
\end{equation}
where the second form uses one-dimensional steady conduction across the strip
\citep{carslaw1947conduction} and is the plane-strain form of the classical result
\citep{boley1960theory}. Two consequences follow. First, $\kth$ does not depend on $d$: the
through-thickness gradient adjusts to cancel the thickness. Second, a strip clamped at both ends
carries the constant restoring moment $D_s\kth$ and does not bow ($w\equiv0$); the bowing
mechanism exists only for pinned ends or for an axially non-uniform $\kth(x)$, which fixes the
support choice A7.

The sign chain is, in order: (i) $\qpp>0$ supplies heat to the outer face (B5); (ii) the temperature
therefore rises towards the outer face ($\Gamma>0$), which is the hot face; (iii) the hot fibres
elongate more, so the strip curves with the outer face convex, $\kth>0$; (iv) with pinned ends this
is an arch whose mid-span moves \emph{away} from the channel, $w>0$; (v) the passage height
$H=H_0+w$ increases. The same chain read backwards shows that a bow \emph{into} the passage
would require heating the fluid-facing face, which is a different model (Section~\ref{sec:limits}).

\subsection{Pinned bowing at negligible axial force}
\label{sec:pinned}
With $w^{\prime\prime}=-\kth$ and $w(0)=w(L)=0$,
\begin{equation}
w(x)=\frac{\kth}{2}\,x(L-x),\qquad \wmax^{\mathrm{lin}}=\frac{\kth L^2}{8}.
\label{eq:wlin}
\end{equation}
This is exact within thin-strip theory when the axial force is negligible, which is the state of a
strip free to expand axially.

\subsection{Thermal membrane force and buckling}
\label{sec:buckling}
With the axial displacement blocked at both ends the mean axial strain vanishes, and integration
across the thickness gives the compressive membrane force per unit depth
\begin{equation}
N_\theta=-\frac{E\,d\,\alpha\bar\theta}{1-\nu},\qquad \Nc=|N_\theta|. \label{eq:ntheta}
\end{equation}
The pinned--pinned Euler load of the wide strip is $\ncr=\pi^2D_s/L^2$ \citep{timoshenko1961stability}.
The wide-strip rigidity is the consistent one because plane strain suppresses the anticlastic
curvature mode. Setting $\Nc=\ncr$ gives the critical mean temperature rise
\begin{equation}
\bar T_{\mathrm{cr}}=\frac{\pi^2d^2}{12\,\alpha(1+\nu)L^2},\qquad
M_b=\frac{\bar T_{\mathrm{cr}}}{\bar\theta}, \label{eq:tcr}
\end{equation}
where the buckling margin $M_b$ is used as a guard. The mean wall rise is estimated as
$\bar\theta=\Delta T_t/2+\qpp/h+\Delta T_b/2$: half the through-wall drop, the interface excess over
the local bulk temperature with the rigid-slit coefficient $h=\mathrm{Nu}_{H1}k_f/D_h$ for one heated
wall \citep{shah1978laminar}, and half the bulk rise $\Delta T_b=\qpp L/(\rho UH_0c_p)$.

\subsection{Beam--column bowing at finite membrane force}
\label{sec:beamcolumn}
A pinned strip carrying the compressive force $\Nc$ couples its deflection to that force at leading
order. Retaining the geometric stiffness of the membrane force and freezing the force at its
blocked-state value, the transverse equilibrium of the strip, with total moment $M=D_s(w^{\prime\prime}+\kth)$
vanishing at the pins, is
\begin{equation}
D_s\,w^{\prime\prime\prime\prime}+\Nc\,w^{\prime\prime}=0,\qquad w(0)=w(L)=0,\quad
w^{\prime\prime}(0)=w^{\prime\prime}(L)=-\kth. \label{eq:beamcol}
\end{equation}
With $\mu^2=\Nc/D_s$ the solution is
\begin{equation}
w(x)=\frac{\kth}{\mu^2}\left[\frac{1-\cos\mu L}{\sin\mu L}\sin\mu x+\cos\mu x-1\right],\qquad
\wmax=\frac{\kth}{\mu^2}\left[\sec\frac{\mu L}{2}-1\right], \label{eq:wbc}
\end{equation}
which reduces to \eqref{eq:wlin} as $\mu L\to0$ and diverges at $\mu L=\pi$, the Euler threshold
\eqref{eq:tcr}: the bowing law and the buckling limit are two limits of one relation. The force
is not in fact frozen. With both ends blocked, the mean axial strain contains the arc-stretching
term $\tfrac12\langle w^{\prime2}\rangle$, so the signed force acquires the correction
\begin{equation}
N_\theta(w)=N_\theta+\frac{Ed}{2(1-\nu^2)}\,\langle w^{\prime2}\rangle,\qquad
\Nc^{\mathrm{sc}}=|N_\theta|-\frac{Ed}{2(1-\nu^2)}\langle w^{\prime2}\rangle, \label{eq:nupdate}
\end{equation}
which relaxes the compression. Because the family \eqref{eq:wbc} holds for every fixed force,
\eqref{eq:nupdate} is a scalar fixed-point problem for $\Nc^{\mathrm{sc}}$.

Two deflections are therefore defined, and they answer different questions. The \emph{frozen-force}
value $\wmax$ of \eqref{eq:wbc} with $\Nc=|N_\theta|$ is the exact solution of the stated
equation \eqref{eq:beamcol}; because it neglects the relaxation it overestimates the deflection and
is a conservative upper bound. The \emph{self-consistent} value solves \eqref{eq:wbc} with
$\Nc^{\mathrm{sc}}$ and is the von K\'arm\'an equilibrium of the reduced strip. It is also what the
coupled algorithm of Section~\ref{sec:algorithm} converges to, since it re-evaluates the membrane force at
every iteration. The self-consistent value is therefore quoted as the model's best estimate, and the
frozen-force value as the upper bound; both are reported. The two differ by \NumExcessFrozen\,\% at
the baseline, \NumExcessW\,\% and \NumExcessS\,\% at the heat fluxes of guards G1 and G2
(Section~\ref{sec:guards}), and \NumExcessB\,\% at $\Nc/\ncr=0.5$, where the frozen-force form should
no longer be used.

\subsection{Hydraulic loading and the thermal-to-hydraulic ratio}
\label{sec:hydraulic}
Integrating plane Poiseuille flow \citep{bird1960transport,shah1978laminar} gives the slit pressure
drop
\begin{equation}
\Delta p=\frac{12\mu UL}{H_0^2}. \label{eq:dp}
\end{equation}
With the outlet datum $p=0$ and the casing cavity at the same reference pressure (A9), the gauge
pressure on the strip falls linearly from $p(0)=p(L)+\Delta p$ at its upstream end to
$p(L)=12\mu UL_o/H_0^2$ at its downstream end, with mean $\bar p=\tfrac12[p(0)+p(L)]$. For pinned ends, a uniform load
$\bar p$ acting in the same sense as the thermal bow deflects the strip by
$w_p=5\bar pL^4/(384D_s)$ at the linear level \citep{timoshenko1970elasticity}; the mild variation of
the load is neglected in this estimate. The ratio of the thermal to the hydraulic deflection of the
\emph{same} strip is
\begin{equation}
\Gamma_p=\frac{w_{\mathrm{th}}}{w_p},\qquad
\Gamma_p^{\mathrm{lin}}=\frac{\wmax^{\mathrm{lin}}}{w_p}=\frac{48\,\kth D_s}{5\,\bar pL^2}, \label{eq:gammap}
\end{equation}
where $\Gamma_p^{\mathrm{lin}}$ uses both deflections at zero membrane force. A second like-for-like
form, $\Gamma_p^{\Nc}$, computes both deflections with the same compressive force $\Nc$: the thermal
deflection from \eqref{eq:wbc} and the pressure deflection from the beam--column equation
$D_sw^{\prime\prime\prime\prime}+\Nc w^{\prime\prime}=\bar p$. Both forms compare two responses of one
structure in one state and are therefore directly comparable. A ratio that places the thermal
deflection \emph{with} the beam--column amplification over a pressure deflection \emph{without} it
mixes two states of the structure and is not used. Because $\kth$ is independent of $d$ while
$D_s\propto d^3$ and $\bar p\propto UH_0^{-2}$, $\Gamma_p$ grows with wall thickness and falls with
velocity and with decreasing channel height.

\subsection{Admissibility guards}
\label{sec:guards}
Five guards bound the regime in which the reduced-order model and the weak form of
Section~\ref{sec:fe} are used. (G1) $\wmax/H_0\le0.5$, so that the deformed fluid mesh keeps positive
Jacobians. (G2) $\sigma_{\max}\le0.2\,\sigma_y$, with $\sigma_{\max}$ the sum of the bending stress
$6|M_{\mathrm{mid}}|/d^2$ at mid-span and the membrane stress $E\alpha\bar\theta/(1-\nu)$.
(G3) $M_b\ge2$. (G4) $\Gamma_p\ge5$, the threshold above which the thermal bow is the dominant
deformation. (G5) $\mathrm{Re}_{D_h}\le100$, $\mathrm{Br}<10^{-2}$ and a Richardson number small enough
for the neglect of buoyancy, with $\mathrm{Ri}=g\beta_f\Delta T^{*}H_0/U^2$ and $\Delta T^{*}=\qpp H_0/k_f$.
